\subsection{Operational observations}
\label{subsec:operational}

Two aspects of running the system for two months do not show up in any
performance metric.

Zone certificate binding solved an administrative problem it was not designed
for. Worker assignments spanning cooperating farms, which would otherwise
need separately maintained access-control configuration, resolved through the
zone attribute already embedded in device and user certificates for
operational reasons. Temporal role expiry removed another class of manual
work, with seasonal credentials lapsing on schedule without administrator
intervention or database editing, each lapse leaving an audit record rather
than a silent deletion.

The enrollment protocol needed three firmware iterations before it survived
multi-hour gateway outages, and the recurring failure mode was time. A sensor
waking with an unsynchronised clock cannot evaluate certificate expiry, and
early firmware either refused to transmit or transmitted under an identity it
could not validate. The fix appears as the final block of
Supplement~S1: acquire time from GPS first with NTP as fallback,
and skip expiry enforcement for a single cycle rather than strand the node
when neither is available. A testbed that never loses its time source cannot
surface this class of fault.

\subsection{Energy-measurement scope}
\label{subsec:energy_detail}

The main article's ``Transmission energy and cryptographic cost'' section
reports only the headline reduction in per-event transmission energy
between single-channel and residue-coded LoRa modes. The full measurement:
per-event energy falls from \EnergySingle\,mJ (SD \EnergySingleSD) to
\EnergyCRT\,mJ (SD \EnergyCRTSD), a reduction of \EnergyReductionPct\%
[\EnergyReductionCILow, \EnergyReductionCIHigh]; per-packet airtime falls
from \AirtimeSingle\,ms to \AirtimeCRT\,ms; and the instrumented packets
average \PayloadSingle\,B and \PayloadCRT\,B respectively. Cryptographic
cost on the sensor is negligible at this duty cycle: the signing pipeline
takes \SignTotal\,\textmu{}s (SD \SignTotalSD), comprising Ed25519 signing
at \SignEdDSA\,\textmu{}s, hardware SHA-256 at \SignSHA\,\textmu{}s, eFuse
key read at \SignEfuse\,\textmu{}s and context management at
\SignContext\,\textmu{}s, or \SignMilli\,ms per wake against the
\DutyCycleMinutes-minute cycle. Three qualifications bound what the energy
reduction describes, given here in full since the main text states only
their conclusion.

First, the instrumented payloads are smaller than a complete authenticated
frame: an Ed25519 signature alone is 64\,B before nonce, timestamp and
identity fields are added, so the measured percentage is a reduction on the
instrumented payload, not on the frame a corrected, fully authenticated
deployment would actually transmit. Second, the energy and airtime samples
come from repeated measurement of one instrumented sensor node on the bench,
so the reported interval reflects measurement repeatability on that node
rather than variation across the fifty deployed devices; a multi-device
energy campaign, ideally over the corrected full frame, would be needed to
generalise the reduction across the fleet. Third, at the deployed
30-minute duty cycle, transmission energy is a small fraction of a node's
daily budget against deep-sleep consumption, so even the full measured
reduction moves estimated battery life by roughly one percent; the value of
the measurement is that it is cleanly attributable to the encoding choice,
not that it materially extends field battery life. We report the reduction
because it is cleanly measured and make no stronger battery-life claim than
that.

\subsection{Signature test vector}
\label{subsec:sig_vector}

A fixed, reproducible vector for the firmware's sign path
(\texttt{esp32/main/signing.c}: SHA-256 the packed struct, then Ed25519-sign
the 32-byte digest), so an independent reimplementation in any language can
check byte-for-byte agreement without running this repository's code.
Field order for the packed struct is little-endian
\texttt{device\_id~(u32)}, \texttt{reading\_id~(u16)},
\texttt{soil\_raw~(u16)}, \texttt{temperature\_centi\_c~(i16)},
\texttt{timestamp~(i64)}; the example uses device id 42, reading id 7, raw
soil 2048, temperature 25.00\,\textdegree{}C (centi-degrees 2500) and
timestamp 1735689600 (2025-01-01T00:00:00Z). The private key is the
Ed25519 seed \texttt{00 01 02 \ldots{} 1e 1f} (32 bytes, byte $i$ has value
$i$), used only to make this vector reproducible and not a key from any
deployed device.

\begin{quote}\small\ttfamily\raggedright
seed:~~~~ 000102030405060708090a0b0c0d0e0f101112131415161718191a1b1c1d1e1f\\
pubkey:~~ 03a107bff3ce10be1d70dd18e74bc09967e4d6309ba50d5f1ddc8664125531b8\\
payload:~ 2a00000007000008c4098085746700000000\\
digest:~~ 29cf4c16f38f0a5bc8cef0da69ed7c9e4d570579fcd2d347a8d063895b49599a\\
sig:~~~~~ 5c59e0f3fef48a23c09ae90c1ac638d179e1058569f2469323f43907b4ab8c4\\
\phantom{sig:~~~~~} d701931da85a762a019cda7174e5a6e418326990f5b51207119e7bd83f124fc08
\end{quote}

Ed25519-verifying \texttt{sig} against \texttt{digest} with \texttt{pubkey}
must succeed; verifying it against the SHA-256 digest of \texttt{payload}
with any single byte flipped must fail. \texttt{gateway/tests/test\_signature\_binding.py}
(\texttt{test\_canonical\_vector\_matches\_firmware\_format} and
\texttt{test\_canonical\_vector\_rejects\_single\_byte\_modification}) checks
both directions and is the executable form of this vector.


